% Mycelium Report Template — COMPREHENSIVE shape
% =============================================================================
% Use when: the report IS the artifact and no separate supplement is expected.
% Full methods, tables, and appendices inline. Long-form draft suitable for a
% paper section, a thesis chapter, or a complete handoff to a collaborator.
%
% Copy to analysis/[name]/reports/[name]-report.tex and replace %%PLACEHOLDER%%
% markers. Numbers and terms must be sourced from .manifest.json; do not
% introduce values that have not been registered there.

\documentclass[11pt,a4paper]{article}

% --- Packages ---
\usepackage[utf8]{inputenc}
\usepackage[T1]{fontenc}
\usepackage{amsmath,amssymb}
\usepackage{graphicx}
\usepackage{booktabs}              % Professional tables
\usepackage{hyperref}
\usepackage[margin=1in]{geometry}
\usepackage{natbib}                % Citations
\usepackage{xcolor}
\usepackage{caption}
\usepackage{subcaption}            % Sub-figures
\usepackage{float}                 % [H] placement
\usepackage{longtable}             % Multi-page tables
\usepackage{listings}              % Code listings

\hypersetup{
    colorlinks=true,
    linkcolor=blue,
    citecolor=blue,
    urlcolor=blue
}

% --- Generated reportable values (see report-values-guide.md / scitexlintr) ---
\IfFileExists{build/report_values.tex}{\input{build/report_values.tex}}{%
  \providecommand{\SciVal}[2]{#1}%
  \providecommand{\SciText}[2]{#1}%
}

% --- Metadata ---
\title{%%TITLE%%}
\author{%%AUTHOR%%}
\date{%%DATE%%}

\begin{document}

% =============================================================================
% TITLE PAGE
% =============================================================================
\maketitle

% =============================================================================
% ABSTRACT
% =============================================================================
\begin{abstract}
%%ABSTRACT%%
% Lead with the scientific question, then the approach in one sentence, then
% the key finding (sourced from the manifest). State the single biggest
% caveat. No undefined acronyms, no citations, no figure references.
% 150–250 words, 1–2 paragraphs.
\end{abstract}

\newpage
\tableofcontents
\newpage

% =============================================================================
% PROBLEM STATEMENT
% =============================================================================
\section{Problem Statement}

%%PROBLEM_STATEMENT%%
% Frame the question in plain English. Define every concept a non-specialist
% would need. Explain why the answer matters, what a good answer would look
% like, and what the comparison baseline is. The comprehensive shape can
% afford 3–5 paragraphs here.

% =============================================================================
% METHODS
% =============================================================================
\section{Methods}

% --- Definitions ---
\subsection{Definitions}

%%DEFINITIONS%%
% Define every mathematical symbol, technical term, and domain-specific
% concept that appears anywhere in the report. Group related definitions
% together. For coined terms that shadow established literature terms, add
% an explicit "this is NOT the standard X" disclaimer at first use, naming
% the specific way it differs (missing factor of 2, not a p-value, etc.).
% Example:
% \textbf{Log2 fold change} ($\log_2 \text{FC}$): The base-2 logarithm of
% the ratio of expression between two conditions. Range: $(-\infty, +\infty)$.

% --- Overview ---
\subsection{Overview}

%%METHODS_OVERVIEW%%
% 2–3 sentences a non-specialist could follow. No tool names, no parameters,
% no version numbers. Even in the comprehensive shape, the Overview stays
% short — its job is to orient the reader before the Technical Detail.

% --- Technical Detail ---
\subsection{Technical Detail}

%%METHODS_DETAIL%%
% Datasets: name, source, size (n samples, n features), preprocessing.
% Software: tool, version, citation for every tool.
% Parameters: value AND rationale for non-default choices. State defaults too
%   (defaults change between versions).
% Statistical methods: test, assumptions checked, correction method, threshold.
% For load-bearing methodological choices, lead with the intuitive
% explanation before the formal definition.

% =============================================================================
% RESULTS
% =============================================================================
\section{Results}

%%RESULTS%%
% Each result follows the pattern:
% 1. State the question being tested.
% 2. Describe what you'd expect under competing hypotheses.
% 3. Report findings with figures and tables — values sourced from the
%    manifest, labels matching label_canonical.
% 4. Conclude in relation to the question.
%
% For each new aggregation analysis type, include one worked example with
% raw inputs traced to the aggregate (prefer inline sparkline-style
% mini-tables over big monolithic figures).
%
% Figure template:
% \begin{figure}[H]
%     \centering
%     \includegraphics[width=0.8\textwidth]{../outputs/figures/FIGURE_NAME.pdf}
%     \caption{[What it shows]. [How to read it]. [Key takeaway].}
%     \label{fig:LABEL}
% \end{figure}

% =============================================================================
% CONCLUSIONS
% =============================================================================
\section{Conclusions}

%%CONCLUSIONS%%
% Map each conclusion back to the problem statement. State caveats explicitly
% — and at a heading level at least as prominent as the claim. Distinguish
% what the data shows from what it suggests. Do not introduce new results.

% =============================================================================
% NEXT STEPS (optional — delete if not applicable)
% =============================================================================
\section{Next Steps}

%%NEXT_STEPS%%
% Only include if there are clear, actionable follow-up analyses. Delete this
% entire section if there is nothing concrete to say.

% =============================================================================
% PROVENANCE
% =============================================================================
\section{Provenance}

\begin{description}
\item[Analysis directory] \texttt{%%ANALYSIS_DIR%%}
\item[Key scripts] ~
  \begin{itemize}
%%SCRIPT_LIST%%
  \end{itemize}
\item[Git commit] \texttt{%%GIT_COMMIT%%}
\item[Analysis period] %%ANALYSIS_PERIOD%%
\item[Analyst] %%ANALYST%%
\item[Report generated] \today
\item[Manifest] \texttt{.manifest.json}
\item[Compile log] \texttt{.compile-log.md}
\item[Software environment] See \texttt{ENVIRONMENTS\_INSTALLATIONS.md}
\end{description}

% =============================================================================
% APPENDIX — INLINE (comprehensive shape)
% =============================================================================
% In the comprehensive shape, the appendix is part of the same document. It
% carries: failure-mode worked examples, parameter sensitivity sweeps, method
% comparisons, subsample stability, extended tables. Reference these from the
% main text Methods or Results.
\appendix

\section{Failure-mode Worked Examples}

%%APPENDIX_FAILURE_MODES%%
% For each new aggregation analysis type whose main-text worked example
% showed a healthy case, include a contrasting failure-mode example here.
% Same format as the main-text example, but the failure mode is the point.

\section{Parameter Sensitivity}

%%APPENDIX_SENSITIVITY%%
% Threshold sweeps, method comparisons, subsample stability.
% Example:
% \subsection{P-value Threshold Sensitivity}
% Figure~\ref{fig:pvalue_sensitivity} shows the number of significant genes
% across adjusted p-value thresholds.
%
% \begin{figure}[H]
%     \centering
%     \includegraphics[width=0.8\textwidth]{../outputs/figures/supplementary/threshold_sensitivity_pvalue.png}
%     \caption{Significant gene count vs. adjusted p-value threshold.
%     Dashed red line: chosen threshold (0.05).}
%     \label{fig:pvalue_sensitivity}
% \end{figure}

\section{Extended Tables}

%%APPENDIX_TABLES%%
% Long tables that would interrupt the main narrative. Use longtable for
% multi-page tables.

% =============================================================================
% REFERENCES
% =============================================================================
\bibliographystyle{plainnat}
\bibliography{references}

\end{document}
